\documentclass[11pt,a4paper]{article}
\usepackage[T1]{fontenc}
\usepackage{lmodern}
\usepackage[margin=28mm]{geometry}
\usepackage{amsmath,amssymb,amsthm,mathtools,mathrsfs}
\usepackage{microtype}
\usepackage{enumitem,booktabs,array,longtable}
\usepackage[dvipsnames]{xcolor}
\definecolor{TheoremLink}{HTML}{1F4E79}
\definecolor{SectionLink}{HTML}{00656D}
\definecolor{EquationLink}{HTML}{643F7A}
\definecolor{CitationLink}{HTML}{2F6B3C}
\definecolor{URLLink}{HTML}{843C39}
\usepackage[colorlinks=true,linkcolor=TheoremLink,citecolor=CitationLink,
 urlcolor=URLLink,bookmarksnumbered=true,pdfencoding=auto]{hyperref}
\hypersetup{pdftitle={An intrinsic non-Archimedean Bedford--Taylor program},
 pdfauthor={},pdfsubject={Research program; 16 November 2026}}
\newcommand{\MathSecref}[1]{{\hypersetup{linkcolor=SectionLink}Section~\ref{#1}}}
\newcommand{\MathEqref}[1]{{\hypersetup{linkcolor=EquationLink}\eqref{#1}}}
\newcommand{\Thmref}[1]{Theorem~\ref{#1}}
\newcommand{\Conjref}[1]{Conjecture~\ref{#1}}
\newcommand{\Questref}[1]{Question~\ref{#1}}
\newcommand{\Probref}[1]{Problem~\ref{#1}}
\newtheorem{theorem}{Theorem}[section]
\newtheorem{proposition}[theorem]{Proposition}
\newtheorem{corollary}[theorem]{Corollary}
\newtheorem{conjecture}[theorem]{Conjecture}
\theoremstyle{definition}
\newtheorem{definition}[theorem]{Definition}
\newtheorem{question}[theorem]{Question}
\newtheorem{problem}[theorem]{Problem}
\newtheorem{example}[theorem]{Example}
\theoremstyle{remark}
\newtheorem{remark}[theorem]{Remark}
\numberwithin{equation}{section}
\newcommand{\R}{\mathbf R}
\newcommand{\Q}{\mathbf Q}
\newcommand{\Z}{\mathbf Z}
\newcommand{\an}{\mathrm{an}}
\newcommand{\D}{\mathrm D}
\newcommand{\dd}{\mathrm d}
\newcommand{\Cb}{\mathcal C_b}
\newcommand{\PSH}{\operatorname{PSH}}
\newcommand{\supp}{\operatorname{supp}}
\newcommand{\tr}{\operatorname{tr}}
\newcommand{\osc}{\operatorname{osc}}
\newcommand{\Sym}{\operatorname{Sym}}
\newcommand{\im}{\operatorname{im}}
\newcommand{\str}{\mathrm{str}}
\newcommand{\ed}{\mathrm{ed}}
\newcommand{\A}{\mathcal A}
\newcommand{\B}{\mathscr B}
\newcommand{\C}{\mathscr C}
\newcommand{\one}{\mathbf 1}
\newcommand{\rel}{\mathrm{rel}}
\newcommand{\res}{\mathrm{res}}
\newcommand{\gr}{\mathrm{gr}}
\newcommand{\trop}{\operatorname{trop}}
\newcommand{\sm}{\mathrm{sm}}
\setlist[enumerate]{leftmargin=*,itemsep=3pt,topsep=5pt}
\setlist[itemize]{leftmargin=*,itemsep=3pt,topsep=5pt}
\setlength{\emergencystretch}{2em}
\setcounter{tocdepth}{2}

\title{An intrinsic non-Archimedean Bedford--Taylor program:\\
 analytic relations, relative capacities and source geometry}
\author{}
\date{16 November 2026}

\begin{document}
\maketitle
\begin{abstract}
The local Bedford--Taylor construction with a positive closed background
provides a common setting for curvature products, relative capacity,
finite energy and measure equations on Berkovich spaces.  We formulate a
program for extending this construction to an intrinsic calculus and for
determining the correct range of its potential theory.  Its principal
conjecture asserts that every actual analytic relation between admissible
finite PL or piecewise-smooth curvature products acts trivially on every allowed
background.  The quadratic problem has a degree-one transgression; it is
settled for finite presentations of residue rank at most two.  On supplied
smooth normal-three model germs, the remaining action is a central
coefficient expressed by a paired spherical tensor and a realized positive
closed tangent current.  We use these results to identify the next
geometric cancellation problem.  Two further questions concern Borel
regularity of relative outer capacity and strong continuity of literal
energy-dual measures.  The nonlinear equation is organized around its
closed source image and the exact normal-probe obstruction.  Explicit
counterexamples determine which capacity, rooftop and solvability
principles can belong to the theory.  We conclude with precise extension
problems for domains, singularities, Green currents and arithmetic
intersection theory.
\end{abstract}
\clearpage

{\hypersetup{linkcolor=SectionLink}\tableofcontents}
\clearpage

\section{Introduction}
\label{sec:introduction}

The Bedford--Taylor operation begins with a simple recursion.  If a
potential can be multiplied by a closed current, its curvature can be
inserted by applying a second derivative.  The force of the construction
comes from proving that the resulting product depends on the curvature,
is symmetric in its differentiated slots, and survives monotone limits.
In the non-Archimedean setting these requirements interact with an
additional geometric issue: a differential expression is represented in
local moment coordinates, and a relation can become zero only after it
is pulled back to the analytic space.  Such a relation must remain zero
when an arbitrary positive closed background is inserted.

The differential forms and currents of Chambert-Loir--Ducros provide the
analytic language for this question \cite{CLD12,CLD25}.  Delta-forms retain
polyhedral boundary terms and connect local calculations with tropical
and arithmetic intersection theory \cite{GK,Mihatsch}.  The global
theory of semipositive metrics and finite-energy potentials provides a
second source of constructions and questions
\cite{BFJmetrics,BFJma,Reboulet}.  The program developed here starts from
the established local calculus with a fixed positive closed background
and asks how far its intrinsic differential, capacitary and variational
structures extend.

Recent tropical work makes the equation side of this distinction
particularly concrete.  Amini--Piquerez--Wu develop regularization,
Bedford--Taylor measures and energy on compact prek\"ahler tropical
varieties, and prove failure of the tropical Calabi--Yau theorem when
a codimension-two star satisfies their genericity condition
\cite[Theorems~5.13 and~11.4]{APW}.  Their comparison with
non-Archimedean solutions uses a strongly faithful tropicalization
and specified support hypotheses \cite[Theorem~12.11]{APW}.  These
results make source compatibility and the analytic realization of
tropical data central questions in the present program.

The choice of background is essential.  A current of positive degree
can detect a small analytic set and can impose additional restrictions
on the sources of a Monge--Amp\`ere operator.  Consequently, the relative
theory contains more than a change of notation in the ambient theory.
Its basic objects are a potential class, an allowed background, a
specified capacity, and, for a global equation, a specified boundary
completion.  We keep these objects explicit throughout.

\subsection{The analytic category and the main conjecture}

Fix a finite extension $K/\Q_p$ with its given absolute value.  An
\emph{analytic space} in this paper is a good, Hausdorff, boundaryless,
quasi-smooth Berkovich $K$-space, locally of pure dimension $n$.  Opens,
germs and sheaves use its ordinary topology.  We write
\[
                         \D=d'd'',\qquad d'd''=-d''d'.
\]
Current identities use the original compactly supported smooth CLD
forms of complementary degree.  A finite moment presentation is refined
by adjoining actual analytic units; every face of the refined image is
retained, including faces on which its tangent map loses rank.

\begin{definition}[Background and classical potential]
\label{def:background}
An \emph{allowed background} on an ordinary open $U$ is a symmetric
CLD middle-positive current $S$ of bidegree $(p,p)$ satisfying
$d'S=d''S=0$.  Middle positivity is the CLD cone detected by the signed
column-square tests $\eta\wedge J\eta$ in complementary degree; its
coefficient matrices are positive semidefinite in those presentations.
Put $q=n-p$.

The class $\Cb(U)$ consists of real-valued, locally bounded,
upper-semicontinuous classical psh functions.  Classical psh means that
pullback to every analytic curve, after every complete isometric
extension of the ground field, is subharmonic.  The bounded
Bedford--Taylor product is denoted by
\begin{equation}
 T_{S,0}=S,\qquad
 T_{S,b}(u_1,\ldots,u_b)
   =\D\bigl(u_1T_{S,b-1}(u_2,\ldots,u_b)\bigr),\quad b\le q.
 \label{eq:bt-recursion}
\end{equation}
For finite differences of such potentials the notation denotes the
signed multilinear extension.
\end{definition}

The existence and properties of \MathEqref{eq:bt-recursion} are
established inputs, recalled in \MathSecref{sec:foundation}.  Let
$\B^{a,b}=B^{a,b}_{\mathrm{Mih}}$ denote the actual analytic delta-form
sheaf, with its full differentials.

\begin{definition}[Admissible finite scalar presentations]
\label{def:admissible}
A finite family of continuous real scalars is \emph{admissible} if,
locally on one ordinary representative, it is expressed through an
actual unit tuple $t_f=(-\log|f_1|,\ldots,-\log|f_N|)$ on its complete
moment image, with one finite coefficient complex, and has the following
properties.  Each scalar is PL or piecewise smooth with a smooth
extension near each closed coefficient cell.  It belongs to the actual
degree-zero delta-form sheaf and has its full curvature
$\D_\delta h\in\B^{1,1}$.  It is also a finite difference of continuous
classical psh functions on that representative.

Write $\mathscr P_{\mathrm{adm}}$ for the resulting local scalar
sheaf.  Its PL subsheaf is denoted by $\mathscr P$.  Every finite
continuous PL presentation in actual unit moments belongs to this
class by local convex-difference decomposition.
\end{definition}

Both the actual delta product and the signed bounded BT product are
therefore defined for an admissible family.  The finite closed-cell
PS extension used in the quadratic results has exactly this scope
\cite{I11,QA}.

For a finite real presentation
\[
 P=\sum_{\nu=1}^{m}c_\nu[h_{\nu1}]\cdots[h_{\nu b}]
       \in\Sym^b_{\R}\mathscr P_{\mathrm{adm}}(U),
\]
define two maps with different targets:
\begin{align}
 \kappa_b(P)
   &=\sum_\nu c_\nu\,
          \D_\delta h_{\nu1}\wedge\cdots\wedge
                             \D_\delta h_{\nu b}\in\B^{b,b}(U),
                    \label{eq:kappa}\\
 \A_S(P)
   &=\sum_\nu c_\nu T_{S,b}(h_{\nu1},\ldots,h_{\nu b}).
                    \label{eq:action}
\end{align}
The brackets in $P$ are formal symmetric-product symbols over $\R$.
They do not denote pointwise multiplication of functions.  We call $P$
an \emph{actual analytic relation} when $\kappa_b(P)=0$ in the analytic
delta-form sheaf.  Its action \MathEqref{eq:action} is an already defined
finite signed coefficient-Radon current.

\begin{conjecture}[Descent of actual analytic relations]
\label{conj:descent}
Let $n\ge1$, $1\le b\le n$, and let $P$ be a finite signed admissible
PL/PS presentation at an ordinary germ $(X,x)$ over the fixed $K$.  Suppose
$\kappa_b(P)=0$ as an actual analytic delta-form germ.  There is an
ordinary representative $U$ of this germ, chosen for $P$, such that
for every $0\le p\le n-b$ and every allowed background $S$ of bidegree
$(p,p)$ on $U$,
\begin{equation}
                              \A_S(P)=0\quad\hbox{on }U.
                              \label{eq:main-descent}
\end{equation}
The equality holds on every original compact smooth complementary CLD
test and after every actual refinement of the presentation.  The field,
its absolute value and the carrier of $S$ remain the given ones.
\end{conjecture}

The quantifiers express the intrinsic content of the conjecture.  The
relation and an ordinary representative are fixed before the background
and its tests are chosen.  A proof would define an action of the actual
curvature-product sheaf $\im\kappa_b$ by choosing any finite local
presentation and using \MathEqref{eq:action}.  Ordinary locality would
then glue these actions.  A common bounded curvature tail can be
absorbed into the background, so the same statement controls its
insertion in larger products.

There are two reasons to begin with $b=2$.  First, this is the first
degree in which actual relations can couple distinct curvature factors.
Second, the quadratic action has a degree-one transgression:
\begin{equation}
 \alpha_P=\frac12\sum_\nu c_\nu
          (h_\nu\D_\delta k_\nu+k_\nu\D_\delta h_\nu),
 \qquad
 \A_S(P)=\D M_S(\alpha_P),
 \label{eq:intro-transgression}
\end{equation}
where $M_S$ is the established degree-one action.  This identity turns
the descent problem into an exactness or controlled-approximation
problem in a module whose action is already available.  In finite
residue rank at most two that construction has been completed
\cite{I11}.  In the next model case, normal dimension three, it reduces
the action to a central angular pairing \cite{I13}.

\subsection{Three interacting directions}

The intrinsic relation problem is one branch of the program.  A second
branch asks which regularity and comparison principles belong to the
relative capacity defined by the full class of normalized classical
testers.  A third studies literal measures on the strong energy
completion and the image of its source map.  These branches share the
bounded calculus, but their remaining arguments have different inputs.
The capacity and source questions can be studied with the existing
potential calculus while relation descent remains open.

This distinction determines the form of the statements below.
\Conjref{conj:descent} is a positive structural conjecture.  Borel
inner regularity of relative outer capacity and literal energy-dual
continuity are posed as questions: their full original formulations
remain undecided.  Nonlinear source compatibility is a classification
problem, because positive smooth density perturbations can already
leave the source image.  The examples in
\MathSecref{sec:capacity} and \MathSecref{sec:sources} are therefore
part of the mathematical specification of the program.

\section{The established potential calculus}
\label{sec:foundation}

We recall the results on which the open problems rest.  Detailed
proofs of the bounded construction, curvature quotient and energy
interfaces are assembled in the companion article \cite{Mature};
the original constructions and more extensive energy theory are in
\cite{R,B,Rchain}.  The statements in this section are used with their
printed domain and boundary hypotheses.

\subsection{Bounded products and the curvature quotient}

The scalar input identifies locally bounded classical psh functions
with limits of decreasing semipositive rational PL sequences on a
single fixed ordinary neighborhood \cite{B}.  In the finite continuous
setting, current positivity is also related to locally uniform smooth
positive approximation in the geometric scope of \cite{R}.  These
regularizations serve different purposes: decreasing approximation
constructs bounded products, whereas uniform approximation is suited
to continuous coefficients and moving continuous data.

For comparison, Wu constructs top-degree local mixed Monge--Amp\`ere
measures for his locally PL-regularizable bounded psh class, with
weighted convergence for decreasing PL approximating rows
\cite[Theorem~4.2, pp.~25--26]{Wu}.  The dependence of a
regularizing neighborhood on the member of a general decreasing psh
sequence is isolated in his Remark~4.6.  His Appendix~A compares the
one-step smooth- and PL-approximable classes; comparison with the
iterated CLD closure uses an envelope property
\cite[Remark~A.2 and Theorem~A.3, pp.~52--53]{Wu}.  The scalar input
used here specifies the classical class over the original $K$ and
one fixed neighborhood; the relative theorem below then addresses
every legal degree and an arbitrary allowed background.

\begin{theorem}[Bounded relative calculus]
\label{thm:bounded}
Let $S$ be an allowed background of degree $(p,p)$ on $U$, and let
$u_1,\ldots,u_b\in\Cb(U)$, where $b\le n-p$.  The recursion
\MathEqref{eq:bt-recursion} defines a coefficient-Radon current,
independent of the decreasing regularizations used in its construction.
It is positive, separately $d'$- and $d''$-closed, symmetric in the
differentiated slots, and additive and positively homogeneous in each
slot.  It commutes with ordinary restriction and gluing.

For decreasing rows $u_{i,j}\downarrow u_i$ in $\Cb(U)$, including a
bounded weight slot $u_{0,j}\downarrow u_0$, one has
\begin{equation}
 u_{0,j}T_{S,b}(u_{1,j},\ldots,u_{b,j})
       \longrightarrow u_0T_{S,b}(u_1,\ldots,u_b)
                         \label{eq:weighted-monotone}
\end{equation}
on all original compact smooth tests.  In top degree this is local
weak convergence of measures on compact continuous scalar tests.
Each differentiated slot depends only on its curvature current.
\end{theorem}

The local CLN bounds control the coefficient measures against a finite
positive frame.  This is the estimate that makes the recursion and
its limiting operations analytic-space constructions.  In
\MathEqref{eq:weighted-monotone} the Borel weight is its given pointwise
representative.  The closure and symmetry assertions concern the
unweighted differentiated product.

Potential independence is stronger than regularization independence.
For bounded classical potentials $u,v$ it gives
\begin{equation}
                  \D[u]=\D[v]\quad\Longrightarrow\quad
                              \D((u-v)S)=0.
                              \label{eq:potential-independence}
\end{equation}
The actual bounded classical curvature kernel is continuous
pluriharmonic.  In the finite continuous class, current harmonicity,
classical pluriharmonicity and ordinary local finite real sums of
logarithms of analytic units agree \cite{R24}.  The same quotient
argument transports curvature order: if $h=a-b$ is a difference of
locally bounded classical potentials, then $\D[h]\ge0$ precisely when
$h$ is classical psh in this class.

Consider the sheaf of actual currents
\begin{equation}
 \mathcal M^1(U)=
   \left\{\sum_{i=1}^m f_i\D[g_i]:
        f_i\in C(U,\R),\ g_i\in\Cb(U),\ m<\infty\right\},
                         \label{eq:curvature-module}
\end{equation}
with signed presentations interpreted by finite expansion.  Products
of elements of this module with a fixed $S$ are already independent
of their continuous-coefficient presentations.  This gives a
substantial intrinsic calculus.  The main conjecture concerns the
additional higher-degree relations arising in the actual delta-form
quotient.

\subsection{Order control and what it measures}

For a positive represented curvature current $Q$, set
\begin{equation}
 E_Q=\{A\in\mathcal M^1:\ -cQ\le A\le cQ
                              \text{ for some }c\ge0\},\qquad
 \|A\|_Q=\inf\{c:\ -cQ\le A\le cQ\}.
                         \label{eq:order-norm}
\end{equation}
Let $\Phi_{S,b}$ denote the established mixed module product.
The positive order estimate is
\begin{align}
 -M_iQ_i\le A_i\le M_iQ_i\quad(1\le i\le b)
 &\quad\Longrightarrow\quad
 -\Bigl(\prod_iM_i\Bigr)R_S
 \le\Phi_{S,b}(A_1,\ldots,A_b)
 \le\Bigl(\prod_iM_i\Bigr)R_S,\nonumber\\[-3pt]
 &\hspace{32mm}R_S=\Phi_{S,b}(Q_1,\ldots,Q_b).
                         \label{eq:order-product}
\end{align}
Its polarized difference estimate is obtained by replacing one factor
at a time.  The completion of $E_Q$ embeds faithfully into actual
coefficient-Radon currents.  The controls $Q$ and $R$ are compatible
through $Q+R$, and ordinary localization gives the sheaf
$\mathcal M^1_{\mathrm{ord}}$.  The mixed products extend uniquely,
with positivity, $C(U)$-multilinearity, locality and support
intersection \cite{Order}.

Order convergence supplies a quantitative comparison with one fixed
positive measure after inserting every background.  Ordinary weak
convergence alone has no such control.  Thus
\MathEqref{eq:order-norm} is useful for relation descent only when the
approximants and the fixed $Q$ are actually constructed.  Its
completion theorem does not supply density of an arbitrary class of
currents.

\subsection{Truncation, no loss and finite energy}

Let $V$ be a strict Laurent box of the following form.  For a compact
original affinoid $B$, choose analytic functions and positive
thresholds such that
\begin{equation}
 \begin{split}
 K_0&=\{|f_a|\le R_a,\ |g_b|\ge r_b\}
                      \Subset\operatorname{Int}_X(B),\\
 V&=\{|f_a|<R_a,\ |g_b|>r_b\}\ne\varnothing,
 \end{split}
                         \label{eq:energy-domain}
\end{equation}
and suppose that $S$ is defined on an ordinary neighborhood of $K_0$.
These assumptions will remain in force for every global energy or
source statement in this paper.  Write $q=n-p\ge1$.

For negative extended classical psh functions $u_1,\ldots,u_b$, put
$u_{i,j}=\max(u_i,-j)$ and $E_j=\bigcap_i\{u_i>-j\}$.  Strict maximum
locality makes
\begin{equation}
 \one_{E_j}T_{S,b}(u_{1,j},\ldots,u_{b,j})
                         \label{eq:nonpolar-truncation}
\end{equation}
a consistent increasing family of positive coefficient measures.
Its locally coefficient-finite limit is the non-polar product.  The
limit is separately closed; closedness is a property of the resulting
current, obtained by truncation estimates.

For a single negative potential $u$, impose
\begin{equation}
 \overline{\{u<-t\}}\Subset V\quad(t>0),
 \qquad
 \nu_{u,j}:=\one_{\{u\le-j\}}T_{S,q}(u_j^{[q]}),
 \qquad \nu_{u,j}(V)\longrightarrow0.
                         \label{eq:no-loss}
\end{equation}
These conditions define $\mathcal F_{S,\partial}(V)$.  If $\mu_u$
denotes the non-polar top measure, define
\begin{equation}
 \mathcal E^\sigma_{S,\partial}(V)
   =\left\{u\in\mathcal F_{S,\partial}(V):
                e_\sigma(u):=\int_V(-u)^\sigma\,\dd\mu_u<\infty\right\},
                  \qquad \sigma>0.
                         \label{eq:energy-class}
\end{equation}
The relative capacity of \MathSecref{sec:capacity} gives exact tail
characterizations:
\begin{align}
 u\in\mathcal F_{S,\partial}(V)
   &\ \Longleftrightarrow\
       t^q C_S(\{u<-t\};V)\longrightarrow0,
                        \label{eq:no-loss-capacity}\\
 u\in\mathcal E^\sigma_{S,\partial}(V)
   &\ \Longleftrightarrow\
       \int_0^\infty t^{q+\sigma-1}
                    C_S(\{u<-t\};V)\,\dd t<\infty,
                        \label{eq:energy-capacity}
\end{align}
within the stated boundary class.  The theory also provides mixed
products in every legal degree, the corresponding recursive current
identities, and mixed integrability estimates.  Uniform energy bounds,
together with the stated capacity convergence, give subcritical
stability; the endpoint requires uniform tail
integrability \cite{Rchain}.

For first energy, the global Green formula makes
\begin{equation}
 B_S(u_0,\ldots,u_q)
    =\int_V(-u_0)\,T_{S,q}(u_1,\ldots,u_q)
                         \label{eq:mixed-energy}
\end{equation}
symmetric in all $q+1$ slots.  Its sharp mixed-energy estimate is
\begin{equation}
 B_S(u_0,\ldots,u_q)
       \le\prod_{i=0}^{q}e_1(u_i)^{1/(q+1)}.
                         \label{eq:mixed-holder}
\end{equation}
The strong completion
$\C=\widehat{\mathcal E}^{1}_{\str}(V)$ has actual Borel
representatives modulo the relative local outer-null ideal, extended
mixed currents $\widehat N_b$, and top measures
$\widehat\mu_u=\widehat N_q(u^{[q]})$.  Put $d=q+1$ and
\begin{equation}
 e(u)=B_S(u^{[d]}),\qquad r(u)=e(u)^{1/d},\qquad
 I(u,v)=e(u)+e(v)-B_S(u,v^{[q]})-B_S(v,u^{[q]}).
                         \label{eq:strong-discrepancy}
\end{equation}
Strong convergence $I(u_j,u)\to0$ is equivalent to convergence in all
local truncated outer capacities together with $e(u_j)\to e(u)$.
The boundary of $\C$ is encoded by strong approximation from
\MathEqref{eq:energy-class}.  Its members need not have ambient
classical representatives.

There are consequently three distinct completions: the order
completion of curvature currents, the strong completion of potentials,
and the Borel completion defined by relative capacity.  Each is useful
because it preserves a different set of operations.  The later
interface problems specify when information passes between them.

\section{Analytic relations and quantitative certificates}
\label{sec:relations}

\subsection{The quadratic transgression}

The degree-one delta action is a canonical map
\begin{equation}
 M_S:\B^{1,1}(U)\longrightarrow\mathcal D'^{p+1,p+1}(U),
 \qquad M_S(f\D_\delta g)=f\D(gS),
                         \label{eq:degree-one-action}
\end{equation}
for the finite continuous scalar presentations at issue.  It is
compatible with actual equality and with the full first derivatives
\cite{Rchain,QA}:
\begin{equation}
 M_S(d'\beta)=d'(\beta\wedge S),\qquad
 M_S(d''\zeta)=d''(\zeta\wedge S).
                         \label{eq:first-derivative-action}
\end{equation}
The polyhedral boundary contributions are included in these formulas.

For $P=\sum c_\nu[h_\nu][k_\nu]$, define
\begin{align}
 u_P&=\sum_\nu c_\nu h_\nu k_\nu,
 &A_{0,P}&=\frac12\D_\delta u_P,\nonumber\\
 Z_P&=\frac12\sum_\nu c_\nu
 (d'_Ph_\nu\wedge d''_Pk_\nu+d'_Pk_\nu\wedge d''_Ph_\nu),
 &\alpha_P&=A_{0,P}-Z_P.
                         \label{eq:quadratic-energy}
\end{align}
Here $d'_P,d''_P$ mean the piecewise derivatives of the continuous
scalar functions.  Full delta Leibniz gives
\begin{equation}
          \D_\delta\alpha_P=\kappa_2(P),\qquad
          \A_S(P)=\D M_S(\alpha_P)=-\D M_S(Z_P).
                         \label{eq:quadratic-return}
\end{equation}
Thus the signed quadratic relation retains the \emph{paired} tensor
$Z_P$.  Treating the two scalar families independently would lose the
cross terms that determine its action.

Let
\begin{equation}
 \mathscr J^{1,1}=d'\B^{0,1}+d''\B^{1,0},\qquad
 \mathscr Q^{1,1}=\B^{1,1}/\mathscr J^{1,1},
                         \label{eq:first-order-quotient}
\end{equation}
where the sums are images in the ordinary sheaf category.  The operator
$\D M_S$ vanishes on $\mathscr J^{1,1}$.  For an actual zero relation
$P$, its class $[\alpha_P]$ lies in the Aeppli quotient
$\ker\D_\delta/\mathscr J^{1,1}$.  For a general nonzero quadratic
curvature, $[\alpha_P]$ instead belongs to the larger quotient
\MathEqref{eq:first-order-quotient}; its derivative is that curvature.

\begin{proposition}[Two sufficient quadratic certificates]
\label{prop:certificates}
Let $P$ be an actual quadratic relation on an ordinary germ.  Either
of the following constructions proves its action vanishes for every
allowed background of degree $p\le n-2$:
\begin{enumerate}
\item an actual local identity
      $\alpha_P=d'\beta+d''\zeta$;
\item full-exact symmetric approximants $a_j=d'\beta_j+d''\zeta_j$,
      one fixed positive represented control $Q$, and
      $\varepsilon_j\downarrow0$, all chosen on an ordinary
      representative before the background and tests, such that
      \begin{equation}
       -\varepsilon_jQ\le[\alpha_P]-[a_j]\le\varepsilon_jQ
                       \label{eq:controlled-exact}
      \end{equation}
      in the established order module, with its all-background
      degree-one action estimate.
\end{enumerate}
\end{proposition}

The first implication follows by applying
\MathEqref{eq:first-derivative-action} and $\D$.  For the second, the
order estimate bounds the difference against each fixed compact test
by $\varepsilon_j$ times a fixed positive control measure; distributional
continuity of $\D$ then applies.  These implications are the proved
transgression and order-continuation results \cite{QA,Order}.  The
construction of the certificates is the remaining mathematical work.

The singular term $A_{0,P}$ has a useful role.  It is already full
first-order exact and can be retained while only the regular energy
$Z_P$ is approximated.  The existing counterexample to a compact
continuous-coefficient smooth-form double primitive for every full
polar \cite{R24} rules out an unnecessary regularity demand.  It leaves
the actual full-exact certificate in
\MathEqref{eq:controlled-exact} available.

\subsection{A stronger carrier seminorm route}

There is a second way to certify all-background descent, based on a
fixed absolute cap carrier.  We state it because it exposes the
difference between an action theorem and a stronger family of
quantitative relation certificates.

Let $\mathcal D(U)$ be the real space of finite continuous locally
bounded classical psh differences and let
$\mathcal P^b(U)=\Sym^b_\R\mathcal D(U)$.  For a closed carrier
$F\subset U$, let $\mathcal Z_F^b(U)$ be generated by the following
specified relations: a factor whose restriction near each point of $F$
agrees with an ordinary local pluriharmonic function; an ordinary
pointwise zero product on all of $F$; and a permitted common ambient
real PL cycle identity together with its actual original pullback.
These are the established carrier relation generators of \cite{R24}.
An arbitrary analytic relation is not included as a generator.

For $L\Subset U$, write $C=F\cap L$ and set
$\|h\|_C=\sup_C|h|$ and
$\osc_Cu=\inf_{a\in\R}\sup_C|u-a|$.  Empty-carrier norms are zero.
Define
\begin{equation}
 \begin{split}
 q^{U,b}_{F,L}(P)=\inf\biggl\{
   \sum_\nu |c_\nu|\,\|h_\nu\|_C
                       \prod_{j=2}^{b}\osc_Cu_{\nu j}:\quad
   P=Z+\sum_\nu c_\nu[h_\nu]\prod_{j=2}^{b}[u_{\nu j}],\\[-2pt]
   Z\in\mathcal Z_F^b(U),\quad h_\nu\in\mathcal D(U),\quad
   u_{\nu j}\in\Cb(U)\cap C(U)\biggr\}.
 \end{split}
 \label{eq:carrier-seminorm}
\end{equation}
Every decomposition in the infimum is finite and is made on the same
ordinary $U$.

For the cap reduction, the carrier is
\begin{equation}
 F=\supp\,T_{[O],c}
       ((\log|f_1|-a_1)_+,\ldots,(\log|f_c|-a_c)_+),
                         \label{eq:absolute-cap}
\end{equation}
where $O$ and $f_i$ are the prescribed ordinary localization and
analytic units.  The relative cap for $S$ is carried by this absolute
set.  The carrier is selected by the cap construction, independently
of cancellation in the signed relation.

\begin{question}[Vanishing of the carrier-CLN certificate]
\label{question:carrier}
Let $P$ be an actual signed PL-zero germ of degree $b\ge2$.  On an
ordinary representative, is
\[
                            q^{U,b}_{F,L}(P)=0
\]
for every required compact stage $L$ and the original absolute cap
carrier \MathEqref{eq:absolute-cap}, with the generator space in
\MathEqref{eq:carrier-seminorm} unchanged?
\end{question}

The established cap and CLN argument sends this zero certificate to
$\A_S(P)=0$ for all allowed $S$ \cite{R24}.  It is therefore a
sufficient route to the PL part of \Conjref{conj:descent}.  The reverse implication
would require information on every polar of the seminorm, whereas
background actions constitute a particular family of functionals.
No identification of these two families is presently available.
Direct exactness in \MathEqref{eq:quadratic-return} may prove the
desired action without settling \Questref{question:carrier}.

There is also a background-approximation route.  Let
$\mathcal W^p_{\mathrm{loc}}$ be the ordinary-local weak closure, on
fixed representatives, of positive separately closed actual
delta-form backgrounds of bidegree $(p,p)$.  The full actual
delta-testing calculus has been established in this class, which is
stable under insertion of bounded BT curvature tails \cite{Rchain}.

\begin{question}[Local approximation of backgrounds]
\label{question:background-density}
Does every allowed background of degree $(p,p)$ belong to
$\mathcal W^p_{\mathrm{loc}}$?  More generally, which intrinsic
geometric or coefficient conditions characterize this local weak
closure on the original analytic space?
\end{question}

A density theorem would transfer the available testing calculus by
its specified weak continuity.  It is an alternative source of
all-background results; the quadratic exactness route treats $S$
directly and can be pursued independently.

\subsection{Higher degrees and differential interfaces}

For $b>2$, a relation can have different curvature tails in its different
summands.  Absorbing one common tail into $S$ treats a useful class of
relations, but leaves this coupled case.  A possible proof architecture
would produce transgressions for coupled presentations and compatibility
relations between them, or prove that their analytic relation ideal is
generated by already controlled lower-degree relations.  Either route
must retain the full boundary differentials on common refinements.

\begin{problem}[From curvature products to differential calculus]
\label{problem:interfaces}
Construct, in the original analytic category, the additional interfaces
needed for a graded differential action: descent for higher coupled
relations, descent with variable finite piecewise-smooth coefficients,
compatibility of the resulting differential operators, and the extension
from ordinary-current tests to the intended delta-form test classes.
For each interface, identify its object sheaf and prove its local
presentation independence, differential identities and stated
functoriality.
\end{problem}

This formulation allows the different interfaces to use different
arguments.  For example, equality on all ordinary smooth tests gives
an ordinary current; extending the pairing to singular delta tests
requires a continuity or integrability theorem for those tests.
Similarly, an analytic map pulls back the scalar potentials used in
a product in its established range, while a pullback of an arbitrary
current requires separate hypotheses.  The resulting calculus should
record the mapping class as part of each naturality statement.

\section{Residue rank and the first critical normal model}
\label{sec:rank}

The finite moment presentation determines which residue geometry enters
the quadratic relation.  This gives a measure of complexity that is
different from the ambient analytic dimension.

\subsection{The completed rank-two construction}

For a fixed actual unit tuple $f=(f_1,\ldots,f_N)$ at $x$, put
\begin{align}
 \mathscr F_f(x)&=\widetilde K^{\gr}
                  (\widetilde f_1(x),\ldots,\widetilde f_N(x)),\nonumber\\
 d_f(x)&=\operatorname{trdeg}^{\gr}_{\widetilde K^{\gr}}
                           \mathscr F_f(x),\nonumber\\
 r_f(x)&=\dim_\Q\operatorname{Span}_\Q
     \{-\log|f_i(x)|+\Q\log|K^\times|\}_{i=1}^N,
 \qquad s_f(x)=d_f(x)-r_f(x).
                         \label{eq:residue-rank}
\end{align}
The number $s_f(x)$ is the transcendence degree of the finitely
generated field of normalized residues of the tuple.  It is bounded by
$s_K(x)=\operatorname{trdeg}_{\widetilde K}\widetilde{\mathcal H(x)}$.

\begin{theorem}[Quadratic descent in finite residue rank at most two]
\label{thm:rank-two}
Fix a finite continuous PL or PS scalar presentation through $f$,
whose coefficient functions extend smoothly near the closed cells of
one finite coefficient complex.  Let
$P=\sum c_\nu[h_\nu][k_\nu]$ be an actual quadratic relation at $x$.
If $s_f(x)\le2$, an ordinary neighborhood and a finite actual
original-unit refinement carry a continuous finite-PS covector whose
full strain is $Z_P$.  In particular $Z_P$, and hence $\alpha_P$,
is full first-order exact there.

There are also full-exact symmetric approximants satisfying
\MathEqref{eq:controlled-exact} with one fixed represented positive
control, chosen before every subsequent background, test and refinement.
Consequently $\A_S(P)=0$ for every allowed $S$ of degree $p\le n-2$.
\end{theorem}

This is the construction of \cite{I11}, not merely a reduction to an
unproved normal-surface lifting statement.  It first realizes the
finite graded moment field on a target over the same $K$.  A complete
normal-surface construction produces the covector on the whole actual
moment image.  Its coefficients then return to the original source;
their strain identities restrict through every later tangent map.
The source-side order estimate retains $A_{0,P}$ and regularizes the
energy with one fixed positive control.  Thus the ambient dimension
and the ranks of later tests are unrestricted.  The rank hypothesis
belongs to the fixed relation tuple.

There is an intrinsic quotient consequence.  Let
$\mathscr K_2=\im_{\mathrm{sh}}\kappa_2$.  At every point with
$s_K(x)\le2$, the formula
\begin{equation}
 \tau_{2,x}:\mathscr K_{2,x}\longrightarrow\mathscr Q^{1,1}_x,
 \qquad \tau_{2,x}(\kappa_2(P))=[\alpha_P]=-[Z_P],
 \qquad \overline\D\,\tau_{2,x}=\operatorname{id}
                         \label{eq:quadratic-quotient}
\end{equation}
is well defined.  Comparing two presentations uses their finite union
tuple; its residue rank is still at most two under this stalk
hypothesis.  On an ordinary open set where $s_K\le2$ everywhere,
these maps form a sheaf morphism.  In particular they do so on analytic
spaces of dimension at most two \cite{I11}.

At a point of larger residue dimension, separate rank bounds on two
presentations do not give a rank bound on their union.  For example,
at the three-dimensional torus Gauss point the pairs $(T_1,T_2)$ and
$(T_2,T_3)$ each generate two residue variables, while their union
generates three.  This explains why the stalk hypothesis in
\MathEqref{eq:quadratic-quotient} has mathematical content.

\subsection{Ordinary punctured descent in a supplied normal-three germ}

The next completed reduction concerns a supplied smooth integral
normal-three model germ $(Y,y)$ over $K^\circ$.  Its distinguished
analytic point has
\begin{equation}
 |\mathcal H(y)^\times|=|K^\times|,
 \qquad F=\widetilde{\mathcal H(y)},
 \qquad\operatorname{trdeg}_{\widetilde K}F=3.
                         \label{eq:normal-three-model}
\end{equation}
Here ``supplied model germ'' includes the original-field lifting and
complete-source comparison of \cite{I13}: the finite residue tuple is
prepared on a smooth projective simple-normal-crossings model of $F$;
its labelled monomial cones, their calibrations and all lower strata
are realized by actual Gauss families on the ordinary germ.  The
comparison covers the complete source of the fixed tuple, and the
families converge ordinarily to $y$ uniformly in their angular
parameters.  These are the model hypotheses of the result below.

\begin{theorem}[Punctured quadratic action and a central residual]
\label{thm:punctured}
Let $(Y,y)$ be a supplied normal-three model germ as above, and let
$P$ be a fixed finite continuous closed-cell PL/PS quadratic
presentation with $\kappa_2(P)=0$.  There is an ordinary representative
$U$ such that, for every allowed background $S$ of degree $p=0$ or $1$,
\begin{equation}
                        \A_S(P)|_{U\setminus\{y\}}=0.
                         \label{eq:punctured}
\end{equation}
For $p=0$ the action vanishes on $U$.  For $p=1$ there is a scalar
$c_S(P)$ such that
\begin{equation}
                              \A_S(P)=c_S(P)\delta_y.
                         \label{eq:central-residual}
\end{equation}
All equalities use the original compact smooth tests and actual
refinements.  The representative and its ordinary cover are fixed for
$P$ before $S$ and its tests are chosen.
\end{theorem}

The proof in \cite{I13} covers a fixed finite old source inventory.
Top-cone interiors have local ambient strain extensions.  Positive
two-faces are treated by proper marked residue curves, simultaneous
Picard completion and a final faithful separator.  The remaining
finite ray families are treated by ordinary products with completed
surface core fibers.  Off the old source the rank-two construction
applies.  Each local primitive covers an ordinary neighborhood of the
actual point, and ordinary current locality yields
\MathEqref{eq:punctured}.

The curve calibration explains why the geometric preparation must keep
its labels.  At a marked residue point $a$, let $b_a=m_av_a$ be the
actual order vector with $v_a$ primitive in the saturated normal
lattice, and let $\delta_a=[\kappa(a):\widetilde K]$.  Its weight and
the slope of a newly adjoined rational function $g$ satisfy
\begin{equation}
                 w_a=\delta_am_a,\qquad
                 w_ah_a(g)=\delta_a\operatorname{ord}_a(g).
                         \label{eq:curve-calibration}
\end{equation}
New zero-old-column marks and satellite pieces can enter after the
Picard functions are adjoined.  Applying the last faithful separator
to the completed tuple is what makes the subsequent common-strain
Stokes argument account for those pieces.

Once the action is supported at $y$, coefficient-Radon structure and
separate closedness eliminate point support below top degree.  In top
degree, a Dirac mass remains possible.  The current problem therefore
has an exact scalar form.

\subsection{Critical tangents and the full spherical pairing}

Choose a protected finite actual moment map
$t:U\to\Omega\subset\R^N$ containing the scalar presentation, with
$t(y)=0$, such that $t^{-1}(L)$ is compact in $U$ for every compact
$L\Subset\Omega$.  For an allowed degree-one background define the
real supercurrent
\begin{equation}
                  \langle T_S,\eta\rangle
                       =\langle S,t^*\eta\rangle,
            \qquad \eta\in\mathcal A_c^{2,2}(\Omega).
                         \label{eq:transpose}
\end{equation}
This uses the entire ordinary inverse image.  It is middle-positive,
symmetric and separately closed, of real bidimension $(2,2)$.
With
\[
        \beta=\sum_i d'z_i\wedge d''z_i,\qquad
        \mu_S=\frac12T_S\wedge\beta^2,
\]
the function $r\mapsto\mu_S(B_r)/r^2$ is nondecreasing.  Thus
\begin{equation}
 \mu_S(B_r)=O(r^2),\qquad
 \nu=\lim_{r\downarrow0}\frac{\mu_S(B_r)}{r^2}<\infty.
                         \label{eq:critical-growth}
\end{equation}
Subsequences of the rescaled currents
$T_{S,\delta}=\delta^2(a_\delta)_*T_S$,
$a_\delta(z)=z/\delta$, have conic positive separately closed limits
$T_0$ with trace mass $\mu_{T_0}(B_R)=\nu R^2$ \cite{I13}.
The real growth calculation is the Lelong monotonicity argument for
positive closed supercurrents \cite[Proposition~3.4 and
Corollary~3.5]{Lagerberg}.  Its application here uses the protected
transpose \MathEqref{eq:transpose}: full compact pullback support
retains both closedness equations.  The later PS pairing and the
comparison with the analytic relation are the additional interfaces
proved in \cite{I13}.

The closed-cell scalar extensions have homogeneous PL first jets
$h_\nu^{[1]},k_\nu^{[1]}$.  Their full paired tensor
\begin{equation}
 Z_0=\frac12\sum_\nu c_\nu
 (d'_Ph_\nu^{[1]}\wedge d''_Pk_\nu^{[1]}
       +d'_Pk_\nu^{[1]}\wedge d''_Ph_\nu^{[1]})
                         \label{eq:leading-tensor}
\end{equation}
satisfies $-Cr\beta\le Z_P-Z_0\le Cr\beta$ on the complete protected
inverse image of $B_r$.  The inequality is preserved by rank-losing
tangent maps.  A cutoff $\rho(z/\delta)$ has second derivatives of
size $O(\delta^{-2})$; combining this with
\MathEqref{eq:critical-growth} shows that replacement of $Z_P$ by
$Z_0$ changes its residual pairing by $O(\delta)$.

\begin{theorem}[Spherical formula for the realized residual]
\label{thm:spherical}
In the setting of \Thmref{thm:punctured}, let $T_0$ be any critical
tangent obtained from the same background and protected presentation.
There is a finite positive semidefinite matrix-valued measure
$\mathsf M$ on $\mathbb S^{N-1}$ such that
\begin{equation}
 \mathsf M\theta=0,\qquad
 \operatorname{div}_{\mathbb S}\mathsf M=0,\qquad
 \sigma=\tr\mathsf M,\qquad
 \sigma(\mathbb S^{N-1})=2\nu.
                         \label{eq:angular-constraints}
\end{equation}
For the actual first jets in \MathEqref{eq:leading-tensor},
\begin{equation}
 \begin{split}
 c_S(P)=\sum_\nu c_\nu\biggl(&
      \int_{\mathbb S^{N-1}}h_\nu^{[1]}k_\nu^{[1]}\,\dd\sigma\\
      &-\int_{\mathbb S^{N-1}}
          d_{\mathbb S}h_\nu^{[1]}\otimes
          d_{\mathbb S}k_\nu^{[1]}:\dd\mathsf M\biggr).
 \end{split}
                         \label{eq:spherical-residual}
\end{equation}
The coefficient-wall values are their canonical tangential PS values.
Equivalently, for every compact smooth cutoff $\rho$ equal to one near
zero,
\begin{equation}
 c_S(P)=-\langle T_0,Z_0\wedge\D\rho\rangle,
 \qquad \D(Z_0\wedge T_0)=-c_S(P)\delta_0.
                         \label{eq:tangent-residual}
\end{equation}
The value is independent of the chosen critical tangent and cutoff.
\end{theorem}

The matrix measure is the angular coefficient measure of the conic
current: its radial coefficient disintegration is $r\,\dd r\,\dd\mathsf
M$.  Positivity controls all coefficients by the trace, and separate
closedness gives the spherical divergence equation.  In fact this
construction classifies the corresponding real positive closed cones.
The analytic cancellation problem concerns the pairs $(Z_0,\mathsf M)$
realized jointly by an original analytic zero relation and its allowed
background.

\begin{conjecture}[Cancellation of the realized critical pairing]
\label{conj:critical}
For every supplied normal-three model germ, every actual quadratic
relation in the finite closed-cell PL/PS class, every allowed
degree-one background, and every common protected presentation, the
right-hand side of \MathEqref{eq:spherical-residual} is zero.
\end{conjecture}

This is the precise remaining special case of quadratic descent in
that model category.  It is proved when $\nu=0$, when $Z_0$ has an
actual full-exact identity on an ordinary neighborhood of the protected
central fiber, or when each scalar first jet is one ambient linear
function.  The last case follows from the angular balance identity
\begin{equation}
          \int_{\mathbb S^{N-1}}\dd\mathsf M
             =\int_{\mathbb S^{N-1}}\theta\theta^t\,\dd\sigma.
                         \label{eq:angular-moment}
\end{equation}
Bounded-potential and the established weakly approximable background
classes provide further known cases \cite{Rchain,I13}.

The next geometric question is a wall-flux cancellation.  Cellwise
strain primitives already exist.  Their differences on overlaps are
affine Killing fields, and smooth interpolation across a wall produces
an $\varepsilon^{-1}$ derivative.  The available normal strip estimate
is $O(\varepsilon)$, so the resulting error can have order one.
Separate closedness controls the compatible PS pairings, but this
critical contribution must be eliminated using the actual zero
relation and its joint source geometry.  A successful argument should
express the wall flux in an actual relation module and then annihilate
it under refinement; an estimate that only treats each open cell
cannot decide \Conjref{conj:critical}.

Two existing examples delimit this approach.  A real five-plane fan
has locally normalized angular primitives and a nonzero global
Killing obstruction, with a quantitative $L^1$ gap; it is $\D$-closed
but not separately $d'$- and $d''$-closed.  It rules out automatic
global patching from those local data.  A separate real coefficient
example has a positive closed plane background and a bounded PS tensor
whose second derivative is a Dirac mass; it reaches the critical
$r^2$ growth.  Neither example is an original analytic zero relation.
Their role is to identify the missing structural input, rather than
to decide the analytic conjecture \cite{I13}.

\section{Relative capacity and the correct Borel question}
\label{sec:capacity}

\subsection{The capacity attached to the full tester class}

For $S$ and $V$ as in \MathEqref{eq:energy-domain}, define on Borel
sets $E\subset V$
\begin{align}
 C_S(E;V)
   &=\sup\left\{\int_E T_{S,q}(u^{[q]}):
                   u\in\Cb(V),\ 0\le u\le1\right\},
                         \label{eq:raw-capacity}\\
 C_S^*(E;V)
   &=\inf\{C_S(G;V): E\subset G\subset V,
                                      \ G\text{ ordinary open}\}.
                         \label{eq:outer-capacity}
\end{align}
The same outer formula applies to arbitrary subsets.  The testers are
defined on the whole specified $V$.  They are actual classical
potentials, so the definition keeps the original analytic constraint.

Relative quasicontinuity, bounded decreasing convergence in capacity,
strict maximum locality, compact-gap comparison and the appropriate
relative quasi-everywhere domination are established.  The ambient
background $[1]$ also has its local polar-set theory and its qualified
exhaustion-domain Choquet theory \cite{Rchain}.  For a general $S$, an
ambient polar set may carry background mass.  Thus relative nullity
must be defined by \MathEqref{eq:outer-capacity}; it cannot be replaced
by ambient pluripolarity.

\begin{question}[Borel inner regularity of relative outer capacity]
\label{question:outer}
For every strict Laurent box \MathEqref{eq:energy-domain}, every
allowed extending background $S$, and every Borel $E\subset V$, is
\begin{equation}
 C_S^*(E;V)\overset{?}{=}\sup\{C_S^*(H;V): H\subset E,\ H\text{ compact}\}
                         \label{eq:borel-inner}
\end{equation}
The question includes the local relatively compact case and its
passage to an exhaustion of $V$.
\end{question}

This is an outer-capacity question.  The stronger equality
$C_S=C_S^*$ is false even for compact sets in the original category.
The crossed-box construction gives a full-tester compact gap and also
excludes the corresponding continuity assertion for the raw capacity
on decreasing compact sets \cite{R24}.  These examples leave
\MathEqref{eq:borel-inner} undecided.

\subsection{Exact examples and a restricted positive result}

There is a second obstruction to transplanting the usual extremal-mass
proof.  Let $X=(\mathbf G_{m,K}^n)^{\an}$ with $n\ge2$,
$\lambda=-\log|\pi|$, $t_i=\log|T_i|$, and
\[
 V=\{|t_i|<\lambda\},\qquad
 \beta=\sum_i d't_i\wedge d''t_i,\qquad S_0=\beta^{n-1}\wedge[X].
\]
Put $\tau=t/\lambda$, $K_r=\tau^{-1}([-r,r]^n)$, and
$A_n=n!2^n\lambda^{n-2}$.  The complete original tester class satisfies
\begin{equation}
                 C_{S_0}(K_r;V)=C_{S_0}^*(K_r;V)
                       =A_n\frac{r^{n-1}}{1-r},\qquad 0<r<1.
                         \label{eq:box-capacity}
\end{equation}
For
$h_r=\max\{(\|\tau\|_\infty-r)/(1-r),0\}-1$,
every bounded actual classical $u$ satisfying
$h_r\le u\le0$ and $u\le-1$ on $\operatorname{Int}K_r$ has
\begin{equation}
                  \int_V\D u\wedge S_0=\frac{A_n}{1-r}.
                         \label{eq:obstacle-mass}
\end{equation}
The ratio of \MathEqref{eq:obstacle-mass} to
\MathEqref{eq:box-capacity} is $r^{1-n}$, which is unbounded as
$r\downarrow0$ \cite{Capacity}.  A total relative extremal mass
therefore cannot calibrate this capacity by a universal constant.
The same original smooth-background setting also supplies failure of
strong subadditivity for outer capacity; ordinary subadditivity
survives.  Strong subadditivity is a possible proof mechanism, rather
than a necessary condition for the desired Borel inner regularity.

There is a useful positive range.  If $B$ is a Borel subset of a
compactly contained moment box and
$\mathcal H^{n-1}(B)<\infty$, then for the full saturated set
$E=\tau^{-1}(B)$,
\begin{equation}
 C_{S_0}^*(E;V)=
       \sup_{H\subset E\ {\rm compact}}C_{S_0}^*(H;V),
                         \label{eq:hausdorff-range}
\end{equation}
and saturated compact subsets suffice.  The proof obtains a uniform
Hausdorff bound for the full tester family; finite-measure inner
regularity then makes the discarded remainder small uniformly in
the testers.  It extends to backgrounds $S'$ for which both $S'$ and
$CS_0-S'$ are allowed positive separately closed currents
\cite{Capacity}.

This suggests a concrete route for larger classes: obtain local
uniform remainder estimates for the original tester measures, and
then pass to Borel sets by an exhaustion that preserves those
estimates.  A countable family of dominating Radon measures, or a
replacement by completed sources, would require a separate theorem.
The finite-Hausdorff argument works because it proves precisely the
uniform remainder bound needed for its specified sets.

\begin{problem}[Detection and tester recovery]
\label{problem:tester}
Find intrinsic conditions on $(V,S)$ and a class of sets under which
normalized actual classical testers recover the intended compact
masses and relative quasi-everywhere statements determine the desired
pointwise representatives.  The conditions should distinguish the
relative null ideal from ambient pluripolarity and remain valid for
the full tester class in \MathEqref{eq:raw-capacity}.
\end{problem}

Such conditions could connect capacity to geometric support or a
verified local approximation of $S$.  They would also supply the
representative information needed in the singularity and boundary
problems below.

\section{Literal measures and the geometry of the source image}
\label{sec:sources}

Fix \MathEqref{eq:energy-domain} and its strong completion $\C$.  The
equation of interest is
\begin{equation}
                              \widehat\mu_u=\mu,\qquad u\in\C.
                         \label{eq:source-equation}
\end{equation}
It has two logically distinct requirements: the prescribed measure
must pair continuously with the strong potential topology, and its
functional must lie in the source image.

\subsection{The energy-difference space and literal continuity}

The completed Borel difference space and its norm are
\begin{equation}
 \mathbb D_S=\C-\C,\qquad
 \|h\|_{\ed}=\inf_{h=a-b}\{r(a)+r(b):a,b\in\C\}.
                         \label{eq:ed-norm}
\end{equation}
This is a Banach space, with $\|u\|_{\ed}=r(u)$ for $u\in\C$.
A positive Radon measure $\mu$ is a \emph{literal energy-dual datum}
if it annihilates the relative null ideal, every member of
$\mathbb D_S$ is literally $\mu$-integrable, and
\begin{equation}
 \ell_\mu(h)=-\int_Vh\,\dd\mu,
 \qquad |\ell_\mu(h)|\le C_\mu\|h\|_{\ed}
                       \quad(h\in\mathbb D_S).
                         \label{eq:literal-ed}
\end{equation}
This definition concerns the actual measure and all completed Borel
differences.  The measure may have infinite total mass on $V$.

\begin{question}[Literal strong continuity]
\label{question:literal}
For every fixed allowed $S$ and strict Laurent box, does
\MathEqref{eq:literal-ed} imply
\begin{equation}
       I(u_j,u)\longrightarrow0\quad\overset{?}{\Longrightarrow}\quad
       \int_V|u_j-u|\,\dd\mu\longrightarrow0
                       \qquad(u_j,u\in\C)
                         \label{eq:literal-strong}
\end{equation}
\end{question}

Source-dominated data and specific capacity-dominated classes already
satisfy this continuity.  The general question asks for the missing
uniform integrability mechanism.  Strong convergence controls the
nonlinear energy discrepancy in \MathEqref{eq:strong-discrepancy};
it need not be convergence in the difference norm
\MathEqref{eq:ed-norm}.  Boundedness of a functional for the latter
norm therefore does not settle \MathEqref{eq:literal-strong}.

Let $P_S(a,b)$ be the greatest member of $\C$ lying below $a$ and $b$
in the relative quasi-everywhere order.  Finite rooftops exist.  One
proved sufficient route to \Questref{question:literal} is the
following fixed-background ordered-rooftop continuity assertion:
\begin{equation}
 \left.
 \begin{gathered}
       w\preceq_{\rel}u,\quad u_j\longrightarrow u
                              \text{ strongly in }\C
 \end{gathered}\right\}
 \quad\Longrightarrow\quad
                  e(P_S(w,u_j))\longrightarrow e(w).
                         \label{eq:ordered-rooftop}
\end{equation}
Another sufficient route is sublinear growth of the energy roots of
finite common minorants along subsequences.  For example, a direct
estimate
\begin{equation}
 e(P_S(u_j:j\in F))\le C\sum_{j\in F}e(u_j)
                         \label{eq:finite-rooftop-budget}
\end{equation}
with $C$ independent of the finite cardinality $|F|$, on the rows
needed by the packing criterion, yields a sublinear root budget
$O(|F|^{1/(q+1)})$.  The established reductions and the exact weaker
subsequence criterion are recorded in \cite{Rchain,Rooftop}.

The coefficient-one binary energy bound is false.  In remaining degree
$q=1$, the universal bound
\begin{equation}
                  e(P_S(u,v))\le2(e(u)+e(v))
                         \label{eq:binary-rooftop}
\end{equation}
has optimal constant $2$ across allowed backgrounds, already for
smooth strongly positive backgrounds and bounded actual obstacles
\cite{Rooftop}.  Repeated use of this binary bound gives a quadratic
worst-case energy budget in the number of obstacles and does not
supply \MathEqref{eq:finite-rooftop-budget}.  The counterexamples also
exclude an energy-only modulus uniform over changing backgrounds;
the fixed-$S$ assertion \MathEqref{eq:ordered-rooftop} remains open.
The constant-one mixed-energy inequality
\MathEqref{eq:mixed-holder} concerns a different expression and
remains available.

\subsection{Variational projection and an exact scalar obstruction}

Define the source functional
\begin{equation}
 J(u)(h)=-\int_Vh\,\dd\widehat\mu_u,
                   \qquad J:\C\longrightarrow\mathbb D_S^*.
                         \label{eq:source-map}
\end{equation}
The image $J(\C)$ is norm closed.  The distance
$\|J(u)-J(v)\|_*$ is complete and induces the strong topology.
Let $\mathbb D_{S,\mathrm s}^*$ be the norm-dual functionals
continuous on every strongly convergent cone row.

\begin{theorem}[Projection and the normal-probe criterion]
\label{thm:normal-probe}
For every $\ell\in\mathbb D_{S,\mathrm s}^*$ there is a unique
$u=\Pi(\ell)\in\C$ maximizing
\[
                        \ell(v)-\frac{e(v)}{q+1}
                                    \quad(v\in\C).
\]
Its normal residual $n=\ell-J(u)$ satisfies
\begin{equation}
                    n(v)\le0\quad(v\in\C),\qquad n(u)=0.
                         \label{eq:normal-cone}
\end{equation}
There is one bounded continuous strictly negative actual classical potential
$\psi_0$, depending only on the box, for which every literal positive
Radon datum $\mu$ with $\ell_\mu\in\mathbb D_{S,\mathrm s}^*$ has
\begin{equation}
 \Xi_{S,\psi_0}(\mu):=-n(\psi_0)\ge0,
 \qquad
 \Xi_{S,\psi_0}(\mu)=0
     \Longleftrightarrow\widehat\mu_{\Pi(\ell_\mu)}=\mu
     \Longleftrightarrow\mu\in\widehat\mu(\C).
                         \label{eq:xi}
\end{equation}
When the equation is soluble, its member of $\C$ is unique.  For every
original compact smooth scalar test $\varphi$ there is a finite
constant $C_\varphi$, independent of $S$ and $\mu$, such that
\begin{equation}
 \left|\int_V\varphi\,
       \dd(\widehat\mu_{\Pi(\ell_\mu)}-\mu)\right|
                     \le C_\varphi\Xi_{S,\psi_0}(\mu).
                         \label{eq:xi-error}
\end{equation}
\end{theorem}

The proof of \cite{R12} constructs $\psi_0$ so that both signs of
each original compact smooth direction are admissible small
perturbations of it.  Testing the normal inequality with these two
perturbations proves \MathEqref{eq:xi-error}.  Literal Radon
uniqueness then identifies the source measures when the defect is
zero.  This makes the remaining equation problem explicit: find
conditions on the given data that eliminate the variational normal.

\begin{problem}[A checkable description of compatible sources]
\label{problem:sources}
For a specified geometry $(V,S)$, characterize substantial classes of
literal positive Radon data $\mu$ for which
$\ell_\mu\in\mathbb D_{S,\mathrm s}^*$ and
$\Xi_{S,\psi_0}(\mu)=0$.  The hypotheses should be verifiable from
$\mu$, the geometry and explicitly supplied reference objects.
Determine the corresponding existence, uniqueness and stability
theorems in $\C$, together with precise obstructions outside these
classes.
\end{problem}

The zero-defect criterion is a complete test after the variational
maximizer is known.  \Probref{problem:sources} asks for input
conditions that make that equality follow.  Assuming the desired
source identity or a sharp norm equality involving the unknown
optimizer would restate the conclusion.

\subsection{Why smooth positive data do not give a universal conjecture}

Tropical non-surjectivity is established in \cite{APW}; the relative
source problem asks for its own geometric constraints on the given
analytic domain and background.  The fixed-source obstruction already
appears in remaining degree one
on a two-dimensional analytic torus \cite{SourceObstruction}.  Use
$V=t^{-1}(( -\lambda,\lambda)^2)$ and
$S=(d't_1\wedge d''t_1+d't_2\wedge d''t_2)\wedge[X]$.
Every completed potential has a canonical convex moment trace $g_u$
with
\begin{equation}
          g_u\in H_0^1(( -\lambda,\lambda)^2),\qquad
                        t_*\widehat\mu_u=\Delta g_u.
                         \label{eq:trace-poisson}
\end{equation}
The trace is defined canonically; arbitrary Borel values at individual
Gauss points are not used.

There are fixed strict reference sources $\mu_w$ and smooth positive
densities $f_j\to1$ uniformly such that
\begin{equation}
                          f_j\mu_w\notin\widehat\mu(\C).
                         \label{eq:density-obstruction}
\end{equation}
The data retain literal strong continuity and uniform source domination.
Their zero-boundary real Poisson solutions fail convexity.  Thus the
linear trace identity \MathEqref{eq:trace-poisson} obstructs the
existence of a source potential in the full prescribed strong
completion.  This is a geometric constraint on the positive potential
cone, and it survives arbitrarily small uniform density perturbations.

A prospective local existence theorem should therefore choose a
compatible data manifold or cone around a known reference $w$.
On a class where the mixed expansion is justified, the linearization
is
\begin{equation}
              L_wh=q\,\D h\wedge(\D w)^{q-1}\wedge S.
                         \label{eq:linearization}
\end{equation}
Three ingredients would make an inverse construction meaningful:
an independently described range of $L_w$ with its boundary and
moment conditions; a right inverse with estimates in spaces embedded
in the required potential completion; and stability of that range
under the mixed nonlinear remainder.  The last condition is substantive.
The available finite-angular-mode constructions impose weighted
harmonic moment conditions on the linear inverse, and a second
nonlinear residual can fail them \cite{C2}.  A right inverse on the
initial finite modes alone does not control the nonlinear iteration.

\subsection{The completed linear equation in its natural space}

Remaining degree one means $q=n-p=1$, in arbitrary ambient dimension.
For the specified affinoid part domains and relative faithful energy
form, the signed measure equation has a developed renormalized theory
\cite{Rchain,Mature}.  Here $H$ is the energy space of that part form,
$T_kv=\max(-k,\min(v,k))$, and $\Gamma_S$ is its energy measure.
A finite signed energy-diffuse Radon datum, meaning that its total
variation vanishes on energy-polar sets, has a
unique canonical zero-boundary renormalized solution $v$.  Its
clippings $T_kv\in H$ belong to the energy space, and its finite
energy-quasicontinuous representative $\widetilde v$ satisfies
\begin{equation}
 \D(\Phi(v)S)=\Phi'(\widetilde v)\mu+
                           \Phi''(\widetilde v)\Gamma_S(v)
                         \label{eq:renormalized-chain}
\end{equation}
for $\Phi\in C^2$ with compactly supported derivative, together with
the radiation condition
$k^{-1}\mathcal E(T_{2k}v-T_kv)\to0$.
There are coarea, level-carrier and total-variation stability results.

This theorem identifies a natural solution space for the linear
measure problem.  Its solution need not be an ambient classical psh
function or a member of the nonlinear source cone $\C$.  The
compatibility obstruction in \MathEqref{eq:density-obstruction} and
the renormalized existence theorem are therefore complementary parts
of the relative theory.

\section{Domains, singularities and geometric applications}
\label{sec:applications}

The Laurent-box hypotheses isolate a domain on which barriers,
truncation estimates and global integration by parts are established.
Extending the theory requires replacing the work performed by each
of those hypotheses.

\subsection{Boundary and representative problems}

\begin{problem}[A domain theory with explicit boundary control]
\label{problem:domains}
Find geometric conditions on an ordinary domain $V$ and an allowed
background $S$ that imply the barrier construction, the global Green
formula and uniform boundary tightness needed for the strong
completion.  Determine when exhaustion by relatively compact
subdomains preserves the no-loss class, the source measures and the
specified boundary condition.
\end{problem}

A useful theorem would state whether the boundary is an ordinary
pointwise value, a capacity boundary value, an energy trace or a
boundary defined by approximation.  It would then give the comparison
between those notions under its geometric assumptions.  A negative
exhaustion potential by itself does not record all the estimates used
in the Laurent-box construction; the missing ingredients include
uniform control of cutoff errors and of weighted mass near the
boundary.

\begin{problem}[Classical representatives of completed potentials]
\label{problem:representatives}
Characterize pairs $(V,S)$ and subclasses of $\C$ for which every
completed Borel class has a canonical ambient classical psh
representative.  Establish the continuity, comparison and boundary
properties of that representative, and identify a relative detection
condition that promotes quasi-everywhere equality to the intended
pointwise equality.
\end{problem}

The curvature quotient and relative capacity provide two possible
inputs.  Curvature determines a bounded potential modulo a
pluriharmonic term, while capacity controls the exceptional sets of
a completed representative.  A representative theorem must connect
these two facts using actual regularization or detection; it cannot
identify their quotient spaces by definition.

\subsection{Singularity types and relative mass}

For negative actual potentials, a local singularity comparison
$u\le v+O(1)$ defines a preorder, and mutual comparison defines a
singularity type.  Non-polar products give a way to ask which masses
and mixed energies are determined by such a type.  In a relative
theory the background is part of this question: it may charge an
ambient polar locus, and a relative outer-null set need not be an
ambient polar set.

\begin{problem}[Relative singularity invariants]
\label{problem:singularities}
For explicit classes of negative actual potentials and allowed
backgrounds, determine which non-polar masses, no-loss conditions and
energy integrability properties are invariant under bounded changes
of singularity.  Construct the associated envelopes and comparison
theorems, specifying their ordinary domains, boundary conditions and
relative null ideals.  Describe how these invariants change under
proper geometric modifications in mapping classes for which the
required current operations are defined.
\end{problem}

The capacity-tail formulas
\MathEqref{eq:no-loss-capacity}--\MathEqref{eq:energy-capacity}
provide a concrete starting point: they turn integrability into
sublevel-set estimates.  A useful first extension would combine these
estimates with a controlled model singularity and prove stability of
the truncation mass under a specified modification.  This would
produce information about singularity types without first requiring
the entire graded differential interface.

\subsection{Green currents and intersection theory}

For a cycle $Z$ of codimension $r$, a Green current has bidegree
$(r-1,r-1)$ and satisfies an identity of the form
\begin{equation}
                              \D g_Z+\delta_Z=\omega_Z,
                         \label{eq:green-current}
\end{equation}
with a prescribed representative $\omega_Z$ of its class.  Existing
delta-form theories construct Green objects and intersection
formulas in substantial geometric ranges \cite{GK,Mihatsch}.
The relative BT program suggests an extension in which the metric
potentials or coefficients have finite energy.

\begin{problem}[Finite-energy Green products]
\label{problem:green}
Construct classes of Green currents satisfying
\MathEqref{eq:green-current} for which mixed products and local
intersection pairings are defined by integrable truncation limits.
Prove independence of the admissible presentations and truncations,
Green symmetry, and agreement with the existing smooth, formal and
delta-form products on their common domains.  State proper-intersection
or replacement hypotheses that control products along the cycles.
\end{problem}

For curvature-generated Green expressions, relation descent would
make the local differential presentation intrinsic.  The energy
theory would provide bounds for the limiting integrals.  Products
with a cycle current and their functoriality would still require
the stated intersection and integrability arguments.  This separates
the algebraic role of \Conjref{conj:descent} from the analytic role
of energy estimates in a proposed Green product.

\subsection{An arithmetic application with local and global conditions}

The continuous theory already connects local current positivity with
metric geometry.  In the geometric scope of \cite{R}, continuous
Chern products glue, enjoy the stated smooth-subvariety projection
formula, and identify the prescribed local mixed measures of the
relevant integrable adelic line bundles.  With the exact
boundary-completion identifications assumed there, finite continuous
local CLD-positive weights give positivity of those prescribed
Yuan--Zhang local measures.  The separate geometric adelic-nefness
input retains its own hypotheses \cite{R,YZ}.

\begin{problem}[Adelic finite-energy intersection]
\label{problem:adelic}
For a specified class of adelic metrized line bundles on a
quasi-projective variety, define local finite-energy intersection
terms that extend the established continuous ones.  Find explicit
summability and boundary conditions under which their sum defines a
global intersection pairing, independent of approximation and
compatible with the product formula.  Prove the resulting variation,
projection and positivity statements under stated geometric
positivity hypotheses.
\end{problem}

The local and global requirements can be tested separately.  At a
non-Archimedean place, the local task is to establish the chosen
mixed-energy integral and its change under replacement of a metric.
Across places, one needs a dominating summable estimate and a
compatible treatment of the quasi-projective boundary.  Local
positivity supplies the sign of a measure; an arithmetic positivity
theorem also needs the global geometric assumptions.  This division
would make a finite-energy extension of adelic intersection theory
verifiable one interface at a time.

\section{A proof architecture for the program}
\label{sec:architecture}

The program has three immediate mathematical targets.  The relation
branch seeks cancellation of the critical pairing
\MathEqref{eq:spherical-residual}, followed by a return from supplied
normal-three models to general ordinary germs and by treatment of
higher coupled relations.  The capacity branch seeks the uniform
tester remainder estimates needed for
\MathEqref{eq:borel-inner}.  The equation branch seeks fixed-background
control sufficient for \MathEqref{eq:literal-strong} and independent
conditions describing the closed source image.

These targets suggest distinct forms of progress.  For the first,
one needs an actual geometric wall-flux identity, an actual full-exact
construction, or a counterexample retaining the complete original
analytic relation.  The real angular classification makes the
residual explicit, but analytic joint realization supplies further
constraints that a successful proof must use.  For the second, a
positive result should exhibit a remainder estimate uniform in the
original normalized testers; a negative result should use those same
testers and ordinary Borel sets.  For the third, an existence theorem
should derive compatibility from data and known reference geometry,
while an obstruction should identify the exact constraint that the
datum violates.

\begin{center}
\small
\begin{tabular}{@{}p{0.23\textwidth}p{0.31\textwidth}p{0.37\textwidth}@{}}
\toprule
Direction & Established input & Missing mathematical step\\
\midrule
Intrinsic relations & Bounded calculus, degree-one action,
quadratic transgression, rank-two construction & Critical joint
cancellation; general-germ return; higher and coefficient relations\\[4pt]
Relative capacity & CLN, quasicontinuity, comparison,
qualified Hausdorff control & Uniform Borel remainder control for
the original testers\\[4pt]
Literal data & Strong completion, energy difference norm,
source and capacity dominated classes & Fixed-background uniform
integrability or the corresponding actual counterexample\\[4pt]
Nonlinear sources & Closed source image, variational projection,
normal-probe criterion & Independent compatible data classes and
stability of their nonlinear constraints\\[4pt]
Geometric applications & Green symmetry, local Chern products,
existing delta and adelic theories & Mapping, singularity,
intersection and global summability interfaces\\
\bottomrule
\end{tabular}
\end{center}

The foundations support these directions without a circular dependence.
The bounded recursion is available before the analytic relation
conjecture.  The capacities and strong energy spaces are built from
that recursion.  Literal continuity would enlarge the data to which
the variational projection applies; it would not remove its normal
residual.  Source compatibility is a further condition.  Conversely,
explicit source classes can be constructed and analyzed without first
settling literal continuity for every energy-dual datum.

A complete theory in this sense would give a coherent account of
objects, operations, limits and compatible equations.  It would
include the nonexistence and strict-inequality phenomena already
described here, together with the hypotheses that replace the false
universal assertions.  The present conjectures and problems specify
the mathematical steps from the established potential calculus to
that theory.

\begin{thebibliography}{99}
\raggedright

\bibitem{CLD12}
A.~Chambert-Loir and A.~Ducros,
\emph{Formes diff\'erentielles r\'eelles et courants sur les espaces de
Berkovich}, \href{https://arxiv.org/abs/1204.6277v1}{arXiv:1204.6277v1},
2012.  Original smooth test sheaf.

\bibitem{CLD25}
A.~Chambert-Loir and A.~Ducros,
\emph{Formes diff\'erentielles r\'eelles et courants sur les espaces de
Berkovich}, \href{https://arxiv.org/abs/1204.6277v2}{arXiv:1204.6277v2},
2025.

\bibitem{GK}
W.~Gubler and K.~K\"unnemann,
\emph{A tropical approach to non-Archimedean Arakelov geometry},
Algebra Number Theory \textbf{11} (2017), no.~1, 77--180.
\href{https://doi.org/10.2140/ant.2017.11.77}{doi:10.2140/ant.2017.11.77}.

\bibitem{Mihatsch}
A.~Mihatsch, \emph{$\delta$-Forms on Lubin--Tate Spaces},
Duke Math. J. \textbf{173} (2024), no.~14, 2809--2928;
\href{https://arxiv.org/abs/2112.10018v4}{arXiv:2112.10018v4}.

\bibitem{BFJmetrics}
S.~Boucksom, C.~Favre and M.~Jonsson,
\emph{Singular semipositive metrics in non-Archimedean geometry},
J. Algebraic Geom. \textbf{25} (2016), 77--139.

\bibitem{BFJma}
S.~Boucksom, C.~Favre and M.~Jonsson,
\emph{Solution to a non-Archimedean Monge--Amp\`ere equation},
J. Amer. Math. Soc. \textbf{28} (2015), 617--667.
\href{https://doi.org/10.1090/S0894-0347-2014-00806-7}{doi:10.1090/S0894-0347-2014-00806-7}.

\bibitem{Reboulet}
R.~Reboulet,
\emph{Plurisubharmonic geodesics in spaces of non-Archimedean metrics
of finite energy}, J. reine angew. Math. \textbf{793} (2022), 59--103.
\href{https://doi.org/10.1515/crelle-2022-0059}{doi:10.1515/crelle-2022-0059}.

\bibitem{YZ}
X.~Yuan and S.-W.~Zhang,
\emph{Adelic line bundles on quasi-projective varieties},
\href{https://arxiv.org/abs/2105.13587v9}{arXiv:2105.13587v9}.

\bibitem{Lagerberg}
A.~Lagerberg, \emph{Super currents and tropical geometry},
Math. Z. \textbf{270} (2012), 1011--1050;
\href{https://arxiv.org/abs/1008.2856}{arXiv:1008.2856}, 2010.
Section~3, especially Proposition~3.4 and Corollary~3.5.

\bibitem{Wu}
L.~Wu, \emph{On the local theory of non-archimedean psh functions},
\href{https://arxiv.org/abs/2609.27610v1}{arXiv:2609.27610v1},
23 September 2026.

\bibitem{APW}
O.~Amini, M.~Piquerez and L.~Wu,
\emph{On solutions to tropical Monge--Amp\`ere equations},
\href{https://arxiv.org/abs/2610.25602v1}{arXiv:2610.25602v1},
19 October 2026.

\bibitem{Mature}
\emph{Bedford--Taylor products with a positive closed
background: curvature quotients, relative energy and source equations},
companion manuscript, 16 November 2026.

\bibitem{R}
\emph{Continuous CLD-positive potentials and mixed currents},
internal manuscript R, version 2.6, 15 November 2026.  Theorems on
regularization, mixed currents and local Chern products.

\bibitem{B}
\emph{Locally bounded classical plurisubharmonic functions and mixed currents}, internal manuscript B, edited version 0.1,
10 November 2026.  Theorems 1.5--1.7 and the specified \`etale
compatibilities.

\bibitem{Rchain}
NA--BT research project, \emph{Relative Bedford--Taylor calculus and finite energy},
internal proof collection, R1--R24, November 2026.  The fixed-background
calculus, potential quotient, unbounded products, R11 energy completion
and relative remaining-degree-one equation.

\bibitem{R12}
NA--BT research project, \emph{Energy duals, source images and a single normal probe},
internal proof collection R12, November 2026.  Scope formulas
M2--M18 and the single-probe normal-residual theorem NP6--NP10.

\bibitem{R24}
NA--BT research project, \emph{Carrier-CLN certificates and the analytic relation
problem}, internal proof collection R24, 14 November 2026.  Carrier
seminorm R24.1, absolute carrier R24.1a, and the harmonic and
primitive results in their stated scopes.

\bibitem{QA}
NA--BT research project, \emph{Degree-one action and the quadratic analytic Aeppli
obstruction}, internal proof module, retained in the November 2026
intrinsic-differential collection, formulas QA1--QA19.

\bibitem{Order}
NA--BT research project, \emph{Intrinsic order continuation}, internal manuscript,
14 November 2026.  Fixed-control mixed estimates and the actual
coefficient-current completion.

\bibitem{I11}
NA--BT research project, \emph{Actual normal stars and ordinary
gluing}, internal intrinsic-differential collection,
Round 11, 15 November 2026.  Ordinary quadratic descent in finite
residue rank at most two and intrinsic quadratic transgression,
modules OR and IQ.

\bibitem{I13}
NA--BT research project, \emph{Normal curve detectors and ordinary refinements},
internal intrinsic-differential collection, Round 13,
16 November 2026.  In particular CC1, punctured action PA, protected
transpose MT, leading energy LE, compatible PS pairing PT and
spherical residual SR.

\bibitem{Capacity}
NA--BT research project, \emph{Structure of the original relative capacity}, internal manuscript,
15 November 2026.  Strong-subadditivity obstruction, exact box
capacity and extremal-mass obstruction, finite-Hausdorff saturated
Borel regularity and background-order extension.

\bibitem{Rooftop}
NA--BT research project, \emph{Sharp binary rooftop energy obstruction in original
smooth backgrounds}, internal manuscript, 15 November 2026, with
the fixed-background continuity question and finite-family criteria.

\bibitem{SourceObstruction}
NA--BT research project, \emph{Fixed all-test strict sources need not have a density
neighborhood of solvability}, internal proof collection R17,
November 2026, formulas FP1--FP9 and the fixed-source theorem.

\bibitem{C2}
NA--BT research project, \emph{Finite-mode source compatibility and nonlinear
moment obstructions}, internal source-compatibility collection,
November 2026.  Weighted harmonic-moment conditions and their
specified finite-mode range.

\end{thebibliography}
\end{document}
